\documentclass[12pt, twocolumn, landscape, a4paper]{article}

\usepackage[margin=1in]{geometry}
\usepackage[utf8]{inputenc}
\usepackage[spanish,es-noshorthands]{babel}
\usepackage{amssymb, amsmath, amsbsy}
\usepackage{tikz}
\usetikzlibrary {arrows.meta} 
\usepackage[pdftex]{hyperref}
\hypersetup{pdftitle={Ejemplos de rectas paralelas},pdfauthor={Luis Carrillo Gutiérrez}}

\pagestyle{empty}

\begin{document}
Hallar los valores de\quad {\Large $x$}\quad e\quad {\Large $y$} : \\

%\begin{itemize} \end{itemize}
\begin{tikzpicture}
	\begin{scope}[Stealth-Stealth]
	\draw (-4, 0) -- (6, 0);
	\draw (-3.75, -2.25) -- (5.25, -2.25);
	\draw (-1.25, -4) -- (3.5, 1.75);
	\end{scope}
	\draw (.25, -.56) node {$(15x + 8)^{\,\circ}$};
	\draw (-1.65, -2.85) node {$(3x + 36)^{\,\circ}$};
	\draw (1.5, -1.75) node {$(3y)^{\,\circ}$};

	\draw (5.5, 0.25) node {$r_{\,1}$};
	\draw (4.75, -1.95) node {$r_{\,2}$};

	\draw (1.25, 0) arc (180:261:.5);
	\draw (-.65, -2.25) arc (180:261:.5);

	\draw (.9, -2.25) arc (0:67:.5);
\end{tikzpicture}

\end{document}